% !PRINTED 360
\weberpage{360}
%<<CONT
and $q$ in the product $\Pi$ runs through the primes,
 distinct from one another, that divide $Q$.

If we denote by
\[
\varrho_1, \varrho_2, \ldots, \varrho_\nu
\]
a complete system of representatives of the numbers coprime to $Q$
in $\Omega$, taken to the modulus $Q$, then in the same way
\weq{6}{O = \varrho_1 O_0 + \varrho_2 O_0 + \cdots + \varrho_\nu O_0}
and
\weq{7}{\nu = \psi(Q) = (O, O_0)}
is the index of the divisor $O_0$ with respect to $O$.

The number $\psi(Q)$ we have determined in general in § 168, (3) of Band II,
 and also for the case in which an ideal takes
the place of $Q$. Since here $N(Q) = Q^2$, we have:
\weq{8}{\psi(Q) = Q^2\, \Pi\left(1 - \frac{1}{N(\mathfrak{q})}\right),}
when $\mathfrak{q}$ runs through the ideal prime divisors of $Q$ that are distinct
from one another.

Now (§ 92)
\[
\begin{aligned}
&\text{1. } N(\mathfrak{q}) = q, \quad \text{when } (\Delta,q) = 0,\\
&\text{2. } N(\mathfrak{q}) = q, \quad \text{when } (\Delta,q) = +1,\\
&\text{3. } N(\mathfrak{q}) = q^2, \quad \text{when } (\Delta,q) = -1.
\end{aligned}
\]
We shall denote these three kinds of primes by $q_1$, $q_2$, $q_3$,
 and observe that every prime $q_2$ is the norm of
two different prime ideals $\mathfrak{q}$. Accordingly there follows from (8)
\weq{9}{\psi(Q) = Q^2\, \Pi\left(1 - \frac{1}{q_1}\right) \Pi\left(1 - \frac{1}{q_2}\right)^2 \Pi\left(1 - \frac{1}{q_3^2}\right).}

Now the index of $O'$ with respect to $O$ is, by Bd.\,II, § 2 (4):
\[
j = \frac{\nu}{\mu} = (O, O') = \frac{(O, O_0)}{(O', O_0)} = \frac{\psi(Q)}{\varphi(Q)}
\]
and
\[
\varphi(Q) = Q\, \Pi\left(1 - \frac{1}{q_1}\right) \Pi\left(1 - \frac{1}{q_2}\right) \Pi\left(1 - \frac{1}{q_3}\right);
\]
consequently there results:
\[
(O, O') = Q\, \Pi\left(1 - \frac{1}{q_2}\right) \Pi\left(1 + \frac{1}{q_3}\right),
\]
for which one may then also put (according to 1., 2., 3.):
\weq{10}{j = (O, O') = Q\, \Pi\left(1 - \frac{(\Delta,q)}{q}\right).}
